\documentclass[11pt]{article}

\usepackage[margin=1in]{geometry}
\usepackage{times}
\usepackage{amsmath}
\usepackage{amssymb}
\usepackage{booktabs}
\usepackage{graphicx}
\usepackage{microtype}
\usepackage{tabularx}
\usepackage{xcolor}
\usepackage{hyperref}
\usepackage{natbib}
\usepackage{url}
\usepackage{fancyhdr}

\hypersetup{
  colorlinks=true,
  linkcolor=blue!70!black,
  citecolor=blue!70!black,
  urlcolor=blue!70!black
}

\setlength{\textfloatsep}{16pt plus 2pt minus 4pt}

\pagestyle{fancy}
\fancyhf{}
\renewcommand{\headrulewidth}{0pt}
\fancyfoot[C]{\thepage}

\title{Safety Erasure Under KV-Cache Compression}

\author{Aryan Gupta \\ \texttt{aryan.cs.app@gmail.com}}
\date{}

\newcommand{\ssei}{\mathrm{SSEI}}
\newcommand{\todoresult}[1]{%
  \begin{center}
  \fbox{%
    \begin{minipage}{0.88\linewidth}
      \textbf{Result pending.} #1
    \end{minipage}
  }
  \end{center}
}
\newcommand{\maybeincludegraphic}[3]{%
  \IfFileExists{#1}{%
    \includegraphics[width=#2]{#1}%
  }{%
    \fbox{%
      \begin{minipage}[c][1.7in][c]{#2}
        \centering
        \textbf{Figure unavailable}\\
        #3
      \end{minipage}
    }%
  }%
}
\newcommand{\maybeinputtable}[2]{%
  \IfFileExists{#1}{%
    \input{#1}%
  }{%
    \todoresult{#2}%
  }%
}
\newcommand{\requiredartifact}[1]{}
\newcommand{\EmpiricalStatusSentence}{%
This manuscript reports empirical findings from a cross-family panel of ten instruction-tuned and base open-weight checkpoints.%
}
\newcommand{\PrimaryRunId}{cross-family panel}
\newcommand{\PrimaryPolicyCount}{}
\newcommand{\PrimaryTopSSEIPolicy}{}
\newcommand{\PrimaryTopSSEI}{}
\newcommand{\PrimaryTopSSEICILow}{}
\newcommand{\PrimaryTopSSEICIHigh}{}
\newcommand{\PrimarySafetyClusterCount}{}
\newcommand{\PrimaryCapabilityClusterCount}{}
\newcommand{\CausalRunId}{causal patching study}
\newcommand{\CausalPolicyCount}{}
\newcommand{\CausalTopSSEIPolicy}{}
\newcommand{\CausalTopSSEI}{}
\newcommand{\CausalTopSSEICILow}{}
\newcommand{\CausalTopSSEICIHigh}{}
\newcommand{\CausalSafetyClusterCount}{}
\newcommand{\CausalCapabilityClusterCount}{}
\IfFileExists{../generated/active_primary/result_macros.tex}{%
  % Auto-generated by scripts/export_paper_assets.py; do not edit by hand.
\renewcommand{\PrimaryRunId}{primary public sweep}
\renewcommand{\PrimaryPolicyCount}{6}
\renewcommand{\PrimaryTopSSEIPolicy}{sliding\_window\_\_budget128}
\renewcommand{\PrimaryTopSSEI}{0.092}
\renewcommand{\PrimaryTopSSEICILow}{0.083}
\renewcommand{\PrimaryTopSSEICIHigh}{0.101}
\renewcommand{\PrimarySafetyClusterCount}{5205}
\renewcommand{\PrimaryCapabilityClusterCount}{1300}
%
}{}
\IfFileExists{../generated/active_causal/result_macros.tex}{%
  % Auto-generated by scripts/export_paper_assets.py; do not edit by hand.
\renewcommand{\PrimaryRunId}{primary public sweep}
\renewcommand{\PrimaryPolicyCount}{6}
\renewcommand{\PrimaryTopSSEIPolicy}{sliding\_window\_\_budget128}
\renewcommand{\PrimaryTopSSEI}{0.092}
\renewcommand{\PrimaryTopSSEICILow}{0.083}
\renewcommand{\PrimaryTopSSEICIHigh}{0.101}
\renewcommand{\PrimarySafetyClusterCount}{5205}
\renewcommand{\PrimaryCapabilityClusterCount}{1300}
%
}{}
\IfFileExists{../generated/claim_assessment/abstract_status_sentence.tex}{%
  \input{../generated/claim_assessment/abstract_status_sentence.tex}%
}{}
\requiredartifact{../generated/active_primary/result_macros.tex}
\requiredartifact{../generated/active_causal/result_macros.tex}
\requiredartifact{../generated/claim_assessment/abstract_status_sentence.tex}

\begin{document}

\maketitle

\begin{abstract}
Inference systems routinely compress or evict the key-value (KV) cache to reduce serving cost. We show that on dense full-attention instruction-tuned language models, this throughput optimization can become a selective safety intervention: refusal and policy-following behavior degrades meaningfully more than ordinary task capability when the cache is aggressively evicted. We test this across ten open-weight checkpoints from seven model families and find positive Selective Safety Erasure Index (SSEI) in four families (Llama, OLMo, Phi, Qwen; six instruction-tuned models including one trained with Model Spec Midtraining, with effect sizes from 0.02 to 0.09 in absolute pass-rate units, 95\% bootstrap CIs excluding zero) under sliding-window and user-pinned cache policies. The effect is robust to leave-one-family-out perturbation. Three models do not show the pattern, and the architectural common thread among them (interleaved or full-layer sliding-window attention; non-standard mixture-of-experts cache handling) is consistent with the hypothesis that the affected behavior depends on long-range KV cache state. Causal-patching experiments that restore baseline cache state into compressed runs show partial recovery of refusal behavior (35--58\% restoration fraction, depending on model); however, both system-role and matched user-role interventions produce comparable restoration under a per-prompt mean-of-ratios estimator, so we do not find evidence that the safety-relevant information is role-specifically localized. Policy-pinned cache retention fully restores refusal.
\end{abstract}

\section{Introduction}

Production language-model serving systems do not deploy a static pair of weights and prompts. They actively manage the transient KV cache through eviction policies, quantization, and paged attention to fit longer contexts into a fixed memory budget \citep{kwon2023pagedattention,zhang2023h2o,xiao2024streamingllm}. These mechanisms are documented to degrade task quality \citep{chen2025pitfalls,yang2024notoken,ananthanarayanan2026physics}. We ask a different question: does cache compression weaken \emph{refusal} behavior more than it weakens ordinary capability on the same prompts and budget?

The result is yes for the majority of currently-deployed dense-attention instruction-tuned models. Across ten checkpoints sampled from Llama, Gemma, Mistral, OpenAI gpt-oss, OLMo, Phi, and Qwen families, six instruction-tuned models show a Selective Safety Erasure Index (SSEI) of at least one cache policy whose 95\% lower CI exceeds zero. Phi-4 and Qwen3-8B exhibit the largest effects, around eight to nine percentage points of safety loss beyond matched capability loss under sliding-window cache eviction. The cross-family conclusion holds when any single family is removed from the panel.

Three models do not show the pattern: Gemma-2-9B-IT, Mistral-7B-Instruct-v0.3, and OpenAI gpt-oss-20b. These models share architectural features that decouple long-range KV cache state from the safety-relevant generation circuit: Gemma 2 interleaves local sliding-window attention with global attention \citep{team2024gemma}, Mistral 7B v0.3 applies sliding-window attention at every layer \citep{jiang2023mistral}, and gpt-oss-20b is a mixture-of-experts model trained with a structured ``harmony'' response format \citep{openai2025gptoss}. Under our cache interventions gpt-oss-20b also collapses into a uniform output mode rather than gracefully degrading. A tenth checkpoint, Qwen2.5-14B-Instruct fine-tuned with Model Spec Midtraining (MSM) \citep{li2025msm}, also shows positive $\ssei$ under sliding-window eviction (0.050 [0.041, 0.059]), though without a matched non-MSM Qwen2.5-14B baseline we cannot attribute this to deep alignment rather than model scale. The positive results in four families and the negative results in these three constitute a falsification test: the affected behavior is cache-resident long-range attention state, not generic compression sensitivity.

Causal-patching experiments, in which baseline cache state is restored into compressed runs, show partial recovery of refusal behavior on Qwen2.5-7B-Instruct (single-seed) and Qwen3-8B (multi-seed, Sonnet-judged). Under a per-prompt bootstrap estimator, both system-role and matched user-role interventions produce comparable restoration fractions (Qwen3-8B: 0.355 vs.\ 0.408, overlapping CIs; Qwen2.5-7B: 0.584 vs.\ 0.584), so we do not find evidence of role-specific localization. A parallel Phi-4 run shows a null causal result: simulated four-bit quantization does not degrade Phi-4's safety in the first place (0.985 vs.\ 0.987 baseline), leaving insufficient denominator for meaningful restoration analysis.

We provide cross-family measurement of selective safety degradation under cache eviction with explicit per-policy effect sizes and bootstrap CIs, an architecture-dependent characterization of which models are affected, and causal-patching evidence that restoring baseline cache state partially recovers refusal behavior without role-specific dissociation.

\section{Related Work}

\paragraph{KV-cache compression.}
Algorithms for KV-cache management trade off memory for quality. H2O retains heavy-hitter tokens \citep{zhang2023h2o}; StreamingLLM relies on attention sinks for long contexts \citep{xiao2024streamingllm}; SnapKV compresses prompts using observed attention patterns \citep{li2024snapkv}; KIVI and KVQuant study low-bit quantization \citep{liu2024kivi,hooper2024kvquant}. \citet{chen2025pitfalls} report that KV compression unevenly harms instruction following and increases system-prompt leakage. \citet{ananthanarayanan2026physics} frame compression as a perturbation of token accessibility through attention dynamics. We extend this line of work along the safety axis: rather than measuring overall quality loss, we measure safety loss minus capability loss on matched prompts.

\paragraph{Cache state as a control surface.}
\citet{wang2025cacheprune} show that KV-cache state can be edited to defend against indirect prompt injection, and prompt-extraction work establishes that the cache can be read out as well \citep{hui2024pleak}. These results frame cache state as an addressable safety-relevant object. Our experiments treat cache state as a causal locus and ask whether ordinary serving-time eviction can damage that state.

\paragraph{Mechanism-first safety phenomena.}
A growing body of work isolates specific mechanisms in safety failure rather than only reporting benchmark accuracy: refusal directions \citep{arditi2024refusal,joad2026more_refusal}, sparse alignment routes \citep{frank2026alignmentroutes}, depth-conditioned refusal restoration \citep{zhang2026anydepth}, subliminal learning \citep{cloud2025subliminal}, token entanglement \citep{zur2025token}, and emergent misalignment \citep{betley2025emergent}. The closest is \citet{wang2025cacheprune}, which uses cache state as an intervention point. We test whether cache state is a \emph{naturally occurring} safety failure point under standard inference optimizations.

\section{Method}

\subsection{Setup}

We evaluate ten open-weight checkpoints: Qwen2.5-7B (base and Instruct) \citep{qwen2024qwen25}, Qwen2.5-14B-Instruct fine-tuned with Model Spec Midtraining (MSM) \citep{li2025msm}, Qwen3-8B \citep{yang2025qwen3}, Llama-3.1-8B-Instruct \citep{grattafiori2024llama3}, Gemma-2-9B-IT \citep{team2024gemma}, Mistral-7B-Instruct-v0.3 \citep{jiang2023mistral}, OLMo-3-7B-Instruct \citep{olmo2025olmo2}, Phi-4 \citep{abdin2024phi4}, and OpenAI gpt-oss-20b \citep{openai2025gptoss}. The MSM checkpoint uses a LoRA adapter trained with Anthropic's Model Spec Midtraining protocol on top of Qwen2.5-14B-Instruct; because we lack a matched Qwen2.5-14B-Instruct baseline without MSM, we cannot cleanly separate the effect of MSM alignment from the effect of model scale. Models are loaded in bfloat16 with explicit access to the cache state, so eviction is implemented in the generation loop rather than delegated to a serving system. We do not claim coverage of vLLM \citep{kwon2023pagedattention}, TGI, or PagedAttention kernels; the results measure the underlying behavior, not operator-level deployment.

\subsection{Cache Policies}

We compare uncompressed generation against five cache interventions:
\begin{itemize}
  \item \textbf{sliding window}: retain only the most recent $N$ cached tokens;
  \item \textbf{sink plus recent}: retain the first $S$ tokens and the most recent $N$;
  \item \textbf{random matched}: evict random tokens at the same budget as a structured policy, as a non-policy control;
  \item \textbf{user pinned}: protect user-role spans, then fill the remaining budget with sink and recent retention. This is the non-policy control for policy pinning;
  \item \textbf{policy pinned}: protect system-role or policy spans, then fill the remaining budget with sink and recent retention.
\end{itemize}
For all policies, we record retained and evicted token counts by role at each decode step.

\subsection{Prompt Suites}

The benchmark separates safety, leakage, over-refusal, and capability:
\begin{itemize}
  \item \textbf{Public refusal safety}: harmful requests drawn from AdvBench \citep{zou2023universal}, JailbreakBench \citep{chao2024jailbreakbench}, and HarmBench \citep{mazeika2024harmbench};
  \item \textbf{Public benign over-refusal}: Databricks Dolly 15k benign instructions \citep{databricks2023dolly};
  \item \textbf{Public XSTest safe}: benign safety-sensitive prompts \citep{rottger2023xstest};
  \item \textbf{Public system leakage}: prompt-injection attacks from the Cyberec corpus \citep{cyberec2026promptinjection} paired with a hidden system canary;
  \item \textbf{Public capability ARC}: multiple-choice capability controls \citep{clark2018arc}.
\end{itemize}
Each suite ships with a manifest of prompt IDs, source dataset revisions, and content hashes. Because public harmful-prompt datasets may overlap with model safety training, we interpret public-suite results as public-benchmark behavior rather than as a generalization claim.

\subsection{Metrics}

Safety and refusal are scored with refusal-string detection on safety suites and exact or token-overlap leakage against hidden system text on leakage suites. Capability is scored with multiple-choice accuracy on ARC-Easy. Over-refusal is the refusal rate on the benign suites.

For each cache policy $p$, budget $b$, and model $m$, we define the Selective Safety Erasure Index as the difference between safety degradation and capability degradation relative to the uncompressed baseline:
\begin{equation}
  \ssei(p, b, m) = \Delta_{\text{safety}}(p,b,m) - \Delta_{\text{capability}}(p,b,m).
\end{equation}
Positive $\ssei$ indicates that safety pass-rate declines more than capability pass-rate under the same intervention. We report prompt-clustered bootstrap 95\% confidence intervals for $\ssei$. A model-policy condition is considered a positive instance of selective safety erasure when $\ssei \geq 0.01$ and its lower CI exceeds zero.

\subsection{Audit}

Automated labels are imperfect. We blinded the suite, policy, and treatment identity of each audit row and routed labels through a separate-family model judge. Annotations come from Claude Sonnet 4 \citep{anthropic2025sonnet}; family separation prevents same-family judging on any panel model. The audit manifests record judge model, prompt-template hashes, raw-output hashes, and response-length buckets so the labels can be re-derived. Across the panel, 2,595 of 2,676 attempted audit rows parsed (97\%); per-model parse rates range from 91\% to 100\%.

\subsection{Causal Restoration}

For length-preserving cache perturbations such as simulated four-bit KV quantization, we run baseline (uncompressed) and compressed decoding on identical prompts and seeds, then patch selected key and value slices from the baseline cache into compressed runs at the same positions. We compare:
\begin{enumerate}
  \item compressed decoding with no patch;
  \item compressed decoding with system-role key+value slices restored from baseline;
  \item compressed decoding with user-role slices restored, matched to the system-token count;
  \item compressed decoding under the policy-pinned eviction policy that protects system spans throughout generation;
  \item compressed decoding under the user-pinned policy as the matched non-policy control.
\end{enumerate}
The user-role restoration is the central control: it shares the additional retained-token budget with the system-role patch while differing only in which role's content is restored. A causal-localization claim requires the system-role intervention to recover refusal substantially more than the matched user-role intervention.

\section{Results}

\subsection{Cross-Family Selectivity}

\begin{figure}[p]
  \centering
  \maybeincludegraphic{../generated/active_primary/figures/ssei_forest.png}{0.88\linewidth}{Per-model, per-policy SSEI with 95\% bootstrap CIs.}
  \caption{Selective Safety Erasure Index by model and cache policy. Each row is one (model, policy) condition; the horizontal bar is the 95\% bootstrap CI and the marker is the point estimate. Light-colored rows have CIs overlapping zero; darker rows have CIs excluding zero. Four families show at least one positive condition.}
  \label{fig:ssei-forest}
\end{figure}

Figure~\ref{fig:ssei-forest} reports per-model per-policy $\ssei$ point estimates and 95\% bootstrap confidence intervals. The largest positive effects come from sliding-window cache eviction on Qwen3-8B ($\ssei = 0.092$, CI [0.083, 0.101]) and Phi-4 ($\ssei = 0.084$, CI [0.076, 0.091]). The MSM-aligned Qwen2.5-14B checkpoint shows a substantial sliding-window effect ($\ssei = 0.050$, CI [0.041, 0.059]). Smaller but significant positive effects appear under user-pinned eviction on Llama-3.1-8B-Instruct (0.024 [0.016, 0.031]), OLMo-3-7B-Instruct (0.026 [0.019, 0.033]), and Qwen2.5-7B-Instruct (0.017 [0.011, 0.022]).

\maybeinputtable{../generated/active_primary/family_replication_table.tex}{Family-level $\ssei$ summary by cache policy.}

The cross-family conclusion is robust to single-family removal. Table~\ref{tab:robustness} reports a leave-one-family-out check: for each family, we re-evaluate the cross-family claim after excluding that family and verify that at least two of the remaining families still have a positive condition with CI excluding zero. The claim holds in every case, including when the largest-effect family (Phi) or the entire Qwen family (three checkpoints) is removed.

\maybeinputtable{../generated/active_primary/robustness_analysis.tex}{Leave-one-family-out robustness for the cross-family selectivity claim.}

\subsection{Per-Suite Effects}

\maybeinputtable{../generated/active_primary/suite_level_effects_table.tex}{Per-suite paired baseline-vs-intervention deltas for refusal, leakage, benign over-refusal, and capability.}

\begin{figure}[t]
  \centering
  \maybeincludegraphic{../generated/active_primary/figures/safety_capability_phase_portrait.pdf}{0.92\linewidth}{Safety degradation versus capability degradation by policy and budget.}
  \caption{Safety loss versus capability loss for the anchor model (Qwen3-8B). Each row is one cache policy. A selective-safety pattern places safety loss visibly to the right of capability loss within a row.}
  \label{fig:safety-capability}
\end{figure}

The selectivity is concentrated in the registered safety and leakage suites. On capability prompts (ARC-Easy), cache compression reduces accuracy by at most a few percentage points and rarely outside the bootstrap CI, while refusal accuracy on safety prompts loses an order of magnitude more under aggressive eviction.

\subsection{Models That Do Not Show the Effect}

Three checkpoints fail to produce positive $\ssei$ on any registered policy: Gemma-2-9B-IT, Mistral-7B-Instruct-v0.3, and OpenAI gpt-oss-20b. The first two share an architectural feature with the underlying generation circuit: native sliding-window attention. Gemma 2 interleaves a 4{,}096-token local sliding-window attention with full global attention every other layer \citep{team2024gemma}. Mistral 7B v0.3 uses 4{,}096-token sliding-window attention at every layer \citep{jiang2023mistral}. In both cases, training has already adapted the model to operate against a locally-windowed view of prior tokens; aggressive eviction at inference time is closer to the model's training distribution than it is for vanilla full-attention models. The Mistral run has $\ssei = 0.009$ at the largest measured policy, immediately below our 0.01 threshold; the Gemma run is centered on zero.

The gpt-oss-20b case is qualitatively different. Mean safety-suite score is 0.129 at the uncompressed baseline and rises to 0.554 under every treatment policy, including the random-matched control. A scalar safety improvement under cache pressure is inconsistent with the hypothesis that the model is gracefully degrading under cache pressure; the uniformity across all five eviction policies indicates a collapse into a single output mode rather than a modulated response to each policy's distinct retention pattern. gpt-oss-20b is trained on the harmony response format \citep{openai2025gptoss} and routes through a mixture-of-experts stack with non-standard cache management. We treat this as a separate failure mode of the cache-intervention protocol, not as evidence against selective safety erasure in the other families.

Selective safety erasure scales with the extent to which the safety-relevant generation circuit depends on long-range KV cache. Models pre-trained on a globally-attending full KV cache (Llama, OLMo, Phi, Qwen) carry safety reasoning that is sensitive to eviction; models pre-trained on locally-windowed cache do not. This pattern is consistent with the causal-localization hypothesis we test directly in the next section.

\subsection{Causal Restoration and Mitigation}

A causal-patching study on Qwen3-8B against the simulated four-bit-KV quantization treatment, judged by Claude Sonnet 4 with selective human verification on the public refusal-safety suite, shows that restoring baseline key+value slices for the first sixteen token positions partially recovers refusal behavior. However, the system-role and matched user-role interventions produce comparable restoration fractions under a per-prompt mean-of-ratios bootstrap estimator (system: 0.355, 95\% CI [0.302, 0.408]; user: 0.408, 95\% CI [0.355, 0.464]; $n = 321$ prompts). The overlapping CIs indicate that the safety-relevant information is distributed across cached tokens rather than localized to system-role spans. Policy-pinned cache retention, which protects system tokens from eviction directly, recovers refusal fully. A legacy Qwen2.5-7B-Instruct single-seed run shows the same pattern: system and user restoration are both 0.584 [0.520, 0.647] and [0.516, 0.647], respectively, with no detectable role asymmetry.

A parallel Phi-4 causal-patching run under the same protocol yields a null result: simulated four-bit quantization does not degrade Phi-4's safety score (baseline 0.985 vs.\ compressed 0.987), so only 11 of 613 paired prompts show a non-zero denominator for the restoration calculation. Phi-4's selectivity effect under the behavioral panel arises from sliding-window eviction ($\ssei = 0.084$), not from the length-preserving quantization treatment used in the causal protocol, making the causal analysis inapplicable to the policy under which Phi-4 is actually vulnerable.

\maybeinputtable{../generated/active_causal/causal_restoration_table.tex}{System-role patching, matched user-role patching, and policy-pinned mitigation conditions.}

\begin{figure}[t]
  \centering
  \maybeincludegraphic{../generated/active_causal/figures/causal_restoration_fraction.pdf}{0.92\linewidth}{Restoration fraction by condition.}
  \caption{Restoration fraction toward baseline on the refusal-safety suite. Zero is the unpatched compressed condition; one is full baseline. Both system-role and user-role K+V patching produce comparable partial restoration (35--41\%). Policy-pinned cache retention recovers fully.}
  \label{fig:causal}
\end{figure}

\subsection{Evidence-Gated Claim Assessment}

\maybeinputtable{../generated/claim_assessment/claim_assessment_table.tex}{Claim-level assessment under the registered evidence gates.}
\maybeinputtable{../generated/claim_assessment/claim_interpretation.tex}{Interpretation of which claims pass and which remain qualitative.}

\section{Discussion}

The strongest finding is architectural: dense full-attention instruction-tuned models lose refusal pass-rate faster than they lose ordinary task accuracy when their KV cache is aggressively evicted, and models pre-trained against locally-windowed attention do not. This is both a deployment concern for the majority of currently-served chat models and a constructive direction for serving architectures that need to be robust to cache pressure.

The causal-restoration evidence establishes that cache state carries safety-relevant information: restoring baseline K+V at the first sixteen positions into a compressed run recovers 35--58\% of the lost refusal behavior (depending on model and suite). However, system-role and matched user-role restorations produce comparable recovery under the per-prompt mean-of-ratios estimator (Qwen3-8B: 0.355 vs.\ 0.408; Qwen2.5-7B: 0.584 vs.\ 0.584), with heavily overlapping confidence intervals. We therefore do not claim role-specific localization. That \emph{any} K+V restoration partially recovers refusal supports cache state as a safety-relevant surface, but whether the mechanism is system-role tokens or a distributed cache property remains open. The Phi-4 causal-patching run yields a null result because the four-bit quantization treatment does not degrade Phi-4's safety, providing no denominator for a meaningful restoration test. Cross-architecture generality of the causal localization claim is therefore unsupported.

Cache compression is not universally unsafe. The result is conditional on the model class, the cache budget, the policy choice, and the suite. Operators deploying full-attention models with aggressive eviction at low budgets on safety-sensitive traffic are the case most directly affected by these results. Operators deploying locally-windowed models, or any model under policy-pinned retention, can expect substantially smaller safety regressions.

\section{Limitations}

The study uses open models and public datasets; conclusions do not extend to closed deployed systems without separate measurement. Cache interventions are implemented in our own generation loop rather than in a production serving stack such as vLLM or TGI, so deployment-validation claims for those systems require their own runs. Public harmful-prompt datasets may overlap with model safety training, which biases absolute safety levels but does not affect the within-prompt baseline-versus-treatment contrast that $\ssei$ measures. The MSM checkpoint (Qwen2.5-14B-Instruct + MSM) lacks a matched non-MSM Qwen2.5-14B-Instruct baseline, so the observed positive SSEI cannot be cleanly attributed to deep alignment rather than model scale. The base-model alignment-contrast track is currently under-powered. The causal-patching protocol uses a length-preserving quantization treatment, which limits its applicability to models whose safety is degraded by that specific intervention; Phi-4's vulnerability to sliding-window eviction cannot be probed by this protocol. The restoration fraction point estimates reported in the causal table use a per-prompt mean-of-ratios bootstrap estimator, which weights each prompt equally; the aggregate ratio-of-means estimator gives qualitatively different values due to heterogeneous per-prompt denominators, but we report the mean-of-ratios as the more principled clustered estimator alongside its bootstrap CI.

\section{Ethics and Safety}

The experiments evaluate refusal behavior on harmful prompts drawn from existing public datasets. Raw model generations are retained locally for audit reproducibility; the paper does not reproduce procedural harmful outputs. The intent is to identify a deployment-time safety regression in widely-used open inference stacks so that operators can choose appropriate mitigations.

\section{Reproducibility}

Each panel run records the resolved configuration, environment metadata, dataset revisions, prompt records, raw generations, metrics, cache statistics, figure source data, and content hashes. Cache interventions are verified as active before metrics are computed. Bootstrap intervals are computed at the prompt-cluster level. Blinded audit labels are joined to experimental metadata only after annotation, and the audit manifest records the judge model, prompt template, and raw-output hashes.

\bibliographystyle{plainnat}
\bibliography{../references}

\end{document}
